\chapter{Implementation}
\label{ch:impl_overview}

The design chapter argued for a particular shape: an external control plane that
adapts an unmodified data plane through deterministic actuation, with
machine-learned decisions kept off the request path and gated behind safety
controls. This chapter shows how we built and deployed that shape, in three parts.
Part~A (this part) gives the deployed-system overview and walks through the data,
control, and actuation surfaces that carry every request and every operator
decision. Part~B continues with the decision-plane services (the anomaly detector,
the forecaster, the routing engine, and the autoscaler) and the models behind
them. Part~C closes with the test, benchmark, and observability machinery used to
validate the whole.

We describe the system as it runs: a single \texttt{docker compose} deployment of
roughly a dozen cooperating services, each a small process with one job. The focus
here is on the modules an operator or integrator touches first, namely the load
balancer and its telemetry pipeline, the policy surface, the sidecar that applies
decisions to NGINX, and the operator interface, together with the slice discipline
used to grow them.

\section{Solution Overview}
\label{sec:ima:overview}

SmartLoad deploys as a set of containers orchestrated by Docker
Compose~\cite{merkel2014docker}. Every component is a service in one
\texttt{docker-compose.yml}, with no hidden host process and no manual wiring
step. Bringing the product up is a single \texttt{docker compose up}, and the same
file is the deployment topology, the dependency graph (through \texttt{depends\_on}
health conditions), and the record of which ports each service listens on.

The services split cleanly into the three planes from the design chapter. The
\emph{data plane} carries client traffic and the telemetry derived from it: the
NGINX load balancer, the log-shipping sidecar that turns the access log into
metrics, and the telemetry-ingest service that stores those metrics. The
\emph{decision plane} reads telemetry and emits events: the anomaly detector, the
forecaster, the routing engine, and the autoscaler. The \emph{control plane} holds
the operating policy and exposes the system to people: the policy manager, the
operator UI, and the actuation sidecar that applies decisions to NGINX. Two pieces
of shared infrastructure, TimescaleDB as the metric store and Redis as the control
bus, sit underneath all three, and the OpenTelemetry Collector bridges the data
plane into the metric store.

\cref{fig:ima:stack} shows the full stack grouped by plane, with each service's
published port. The solid path carries client traffic and its derived telemetry;
the control bus then carries decisions back to the actuation sidecar, the only
writer of the NGINX upstream configuration.

\begin{figure}[tbp]
\centering
\includegraphics[width=\linewidth,height=0.88\textheight,keepaspectratio]{system_design}
\caption[The full stack grouped by plane]{The full stack grouped by plane: the SmartLoad services together with the infrastructure containers, with ports matching docker-compose.}
\label{fig:ima:stack}
\end{figure}

\cref{tab:ima:services} names each module, its plane, its implementation language,
and its single job. The main language is Python: every decision-plane service, the
policy manager, the telemetry ingest, the log shipper, and the operator-UI backend
are Python processes. Working in one language across most of the stack lets us
unit-test the pure decision logic of each service without its I/O dependencies. The
data plane itself is NGINX (configuration, not code); the operator UI front end is
TypeScript compiled by Vite; and the backend pool used for load testing is a small
Node.js service.

\begin{table}[H]
\centering
\caption[The deployed services by plane and language]{The deployed services, by plane, language, and responsibility.}
\label{tab:ima:services}
\small
\begin{tabularx}{\linewidth}{@{}l l l Y@{}}
\toprule
\textbf{Service} & \textbf{Plane} & \textbf{Language} & \textbf{Responsibility} \\
\midrule
load-balancer    & Data    & NGINX conf & Terminates client traffic, proxies to the backend pool \\
lb-otel-shipper  & Data    & Python     & Tails the NGINX access log, ships per-request metrics as OTLP \\
resource-collector & Data  & Python     & Ships per-container CPU and memory telemetry to the metrics store \\
otel-collector   & Data    & (collector)& Receives OTLP and forwards to the telemetry service \\
telemetry        & Data    & Python     & Parses OTLP, writes the \texttt{metrics} hypertable, serves a read API \\
test-backend     & Data    & Node.js    & The backend pool being routed across (load-test target) \\
anomaly-detector & Decision& Python     & Classifies each backend healthy / degraded / unhealthy \\
forecasting      & Decision& Python     & Predicts short-horizon request rate for the autoscaler \\
rl-engine        & Decision& Python     & Ranks backends for routing via the deployed monotone heuristic (shadow or active) \\
autoscaler       & Decision& Python     & Grows / shrinks the backend pool via the Docker daemon \\
policy-manager   & Control & Python     & Owns the operating policy, audits changes, publishes them \\
lb-sidecar       & Control & Python     & Applies routing / exclusion / scale decisions to NGINX \\
operator-ui      & Control & Python/TS  & Operator surface: health, policy, audit, actions, live feed \\
\midrule
TimescaleDB      & Shared  & PostgreSQL & Time-series store for metrics and audit history \\
Redis            & Shared  & (server)   & Pub/sub control bus between services \\
\bottomrule
\end{tabularx}
\end{table}

Two design properties of this stack recur through the rest of the chapter. First,
\emph{the decision plane never sits in the request path}. A failure in the anomaly
detector or the forecaster does not stop traffic; it stops \emph{adaptation}, and
NGINX keeps serving under its last applied configuration. Second, \emph{every
cross-service interaction flows through one of exactly two contracts}: an HTTP REST
surface under \texttt{/api/v1} for operators and integrators, and a Redis pub/sub
surface for inter-service events. Both contracts are defined once, in a shared
in-tree module that every service imports rather than redefining the wire format.
That single source of truth keeps a twelve-service deployment from drifting into a
dozen incompatible message formats.

\subsection{Third-party infrastructure}
\label{sec:ima:thirdparty}

SmartLoad's own services run on established, vendor-neutral infrastructure rather
than bespoke components, which lets the whole stack stand up from a single
composition on any Docker host. \cref{tab:ima:thirdparty} lists the external images
the deployment pulls, the pinned version of each, and the role it plays. The
choices are deliberately conservative: a PostgreSQL-derived time-series store, a
single-node Redis for the control bus, the reference OpenTelemetry Collector, and
the standard Prometheus and Grafana pair for service-health observability. The data
plane adds one more, NGINX, built into a custom image but configured rather than
forked. The machine-learning libraries the decision engines depend on
(scikit-learn, statsmodels, Stable-Baselines3, and PyTorch) are inventoried with
their justification in \cref{ch:impl_overview}, Part~B (\cref{tab:imb:tooling});
they are runtime dependencies of individual services, not stack-level
infrastructure.

\begin{table}[H]
\centering
\caption[Third-party infrastructure dependencies]{Third-party infrastructure the deployment depends on, with the pinned image and the role each plays.}
\label{tab:ima:thirdparty}
\small
\begin{tabularx}{\linewidth}{@{}l l Y@{}}
\toprule
\textbf{Component} & \textbf{Image / version} & \textbf{Role in the stack} \\
\midrule
TimescaleDB    & \texttt{timescale/timescaledb:latest-pg16} & Hypertable store for per-request metrics and audit history \\
Redis          & \texttt{redis:7-alpine}                    & Pub/sub control bus for inter-service decisions \\
redis-exporter & \texttt{oliver006/redis\_exporter:v1.62.0-alpine} & Exposes control-bus health to Prometheus \\
OTel Collector & \texttt{otel/opentelemetry-collector-contrib} & Receives OTLP from the log shipper, forwards to telemetry \\
Prometheus     & \texttt{prom/prometheus}                   & Scrapes the per-service \texttt{/metrics} surfaces and exporters \\
Grafana        & \texttt{grafana/grafana}                   & Renders the operational overview dashboard (others planned) \\
NGINX          & custom build over \texttt{nginx}           & The data-plane load balancer (configured, not forked) \\
Docker Compose & (engine)                                   & Single-host orchestration of the whole stack \\
\bottomrule
\end{tabularx}
\end{table}

\subsection{External surfaces and the port map}
\label{sec:ima:ports}

Because the deployment is a single composition, the ports it publishes are also the
complete list of ways the outside world reaches it. \cref{tab:ima:ports} is that
map. It separates the three surfaces an operator or integrator uses directly, the
ingress on \texttt{:8080}, the operator UI on \texttt{:8090}, and the demo console
on \texttt{:8091}, from the observability and tooling surfaces (Grafana,
Prometheus, Locust) and from the per-service \texttt{/api/v1} ports the operator UI
and SDK consume. Only the ingress carries client traffic; every other surface is a
control, observability, or diagnostic plane that can fail without affecting a
single request.

\begin{table}[H]
\centering
\caption[Published external surfaces of the deployment]{Published external surfaces of the deployment, where only \texttt{:8080} carries client traffic.}
\label{tab:ima:ports}
\small
\begin{tabularx}{\linewidth}{@{}l l Y@{}}
\toprule
\textbf{Port} & \textbf{Service} & \textbf{Surface} \\
\midrule
\texttt{:8080} & load-balancer    & Client ingress (NGINX, the only request path) \\
\texttt{:8090} & operator-ui      & Operator transparency and override UI \\
\texttt{:8091} & demo-ui          & Developer ``Dev Console'' (benchmarks, one-click load) \\
\texttt{:8089} & traffic-simulator& Locust load-generation UI \\
\texttt{:3000} & grafana          & Operational dashboards \\
\texttt{:9090} & prometheus       & Metric query and scrape target browser \\
\texttt{:8081} & telemetry        & OTLP ingest and read API \\
\texttt{:8082} & anomaly-detector & Decision-engine \texttt{/api/v1} and \texttt{/metrics} \\
\texttt{:8083} & forecasting      & Decision-engine \texttt{/api/v1} and \texttt{/metrics} \\
\texttt{:8084} & rl-engine        & Routing-engine \texttt{/api/v1} and \texttt{/metrics} \\
\texttt{:8085} & autoscaler       & Autoscaler \texttt{/api/v1} and \texttt{/metrics} \\
\texttt{:8086} & policy-manager   & Policy \texttt{/api/v1} and audit reads \\
\texttt{:8087} & lb-sidecar       & Sidecar state and \texttt{/metrics} \\
\bottomrule
\end{tabularx}
\end{table}

\section{Data Plane and Telemetry Pipeline}
\label{sec:ima:dataplane}

The data plane answers one question end to end: how does a client request become a
row of telemetry the decision plane can reason over? The path crosses five
services, each link chosen to preserve \emph{per-request fidelity}, the property
that the decision plane sees one record per request, not a pre-aggregated summary.

\subsection{NGINX as the data-plane entry point}
\label{sec:ima:nginx}

The load balancer is a stock NGINX~\cite{nginx} container that terminates client
traffic and proxies to a pool of backend replicas through a single
\texttt{upstream} block. The pool is listed explicitly rather than resolved
dynamically: each \texttt{server} line names one backend
(\texttt{smartload-test-backend-1} through \texttt{-5}) with the health directives
\texttt{max\_fails=3 fail\_timeout=10s}, and the proxy is wired with
\texttt{proxy\_next\_upstream error timeout http\_502 http\_503 http\_504}.

Two directives carry the resilience contract. \texttt{max\_fails=3
fail\_timeout=10s} takes a backend out of rotation after three failures inside a
ten-second window, a deterministic, built-in circuit breaker that needs no help
from the decision plane. \texttt{proxy\_next\_upstream} makes failure graceful: if
the chosen backend returns a 5xx or times out, NGINX retries the next backend
instead of showing the error to the client. This matters because the autoscaler
toggles backends between running and stopped at runtime; without it, every scale-in
would reach clients as a transient 502. The five listed slots are a seed, not a
ceiling: the sidecar of \cref{sec:ima:sidecar} overwrites this block whenever a
decision arrives, and the autoscaler can grow the pool beyond five. The static
block is the bootstrap fallback that lets NGINX start before any decision is
published.

One subtlety matters for the evaluation. The \texttt{max\_fails=3} circuit breaker
above belongs to the static bootstrap block; the upstream the sidecar renders once it
takes over, and the configuration the routing benchmark runs, instead set
\texttt{max\_fails=0}, deliberately disabling NGINX's passive ejection. Backend
exclusion is the decision plane's job, not the data plane's, and ceding it cleanly is
what makes the comparison in \cref{ch:impl_testing} meaningful: with \texttt{max\_fails=0}
a backend that is slow or shedding \texttt{503} but never trips the passive check is
\emph{un-ejectable} by plain round-robin, which keeps routing a fifth of its traffic
onto it, while SmartLoad's detector finds it and the sidecar pulls it out. That is the
exact failure mode the \texttt{B\_degrade} and \texttt{D\_slow} phases isolate.

A second, small detail in the NGINX configuration is a static \texttt{location =
/health} that returns SmartLoad's canonical \texttt{\{"status":"ok"\}} for the load
balancer itself, with its own access log suppressed. The exact-match \texttt{=}
intercepts only \texttt{/health} while every other path still proxies to the pool,
keeping the operator UI's health probe from being round-robined into a backend and
misreporting the load balancer's state.

\subsection{From access log to OTLP}
\label{sec:ima:shipper}

NGINX emits one JSON object per request in a custom \texttt{smartload\_json} log
format, carrying the timestamp, status, the upstream that served the request, and
the request latency. Rather than use NGINX's native metrics module, which
aggregates into histograms before the database, SmartLoad ships the \emph{log}
itself, for fidelity: the anomaly and routing queries compute standard deviations
and maxima over rolling windows, and those statistics need per-request rows, not
pre-aggregated histograms. The log-based path is the only one that keeps the
individual record.

The \texttt{lb-otel-shipper} is a Python sidecar that runs alongside NGINX, tails
the access log on a shared volume, and emits one OpenTelemetry data
point~\cite{opentelemetry} per request per metric over OTLP/HTTP-JSON. Each log
line becomes three canonical metrics, \texttt{request\_count},
\texttt{request\_latency\_ms}, and \texttt{error\_rate}, the last computed by the
shipper because NGINX does not emit it directly (any status $\geq 500$ is an
error). The metric names are the contract: the decision-plane SQL filters on
exactly these strings. The shipper batches points and flushes every second or every
five hundred points, whichever comes first, and is strictly fire-and-forget: if the
collector is unreachable, it counts and drops the batch but never blocks the tail
loop. Back-pressure must never reach the data plane.

\subsection{Collector, ingest, and persistence}
\label{sec:ima:telemetry}

The OpenTelemetry Collector receives the shipper's OTLP and forwards it to the
\texttt{telemetry} service, a Flask process that owns the schema contract for the
\texttt{metrics} hypertable in TimescaleDB~\cite{timescaledb}. It is the only
writer of that table, so a schema change touches one file. It parses the four-level
OTLP envelope (\texttt{resourceMetrics} $\to$ \texttt{scopeMetrics} $\to$
\texttt{metrics} $\to$ \texttt{dataPoints}), drops anything that is not a gauge or
sum (histograms and summaries are counted, not silently discarded), and inserts the
surviving rows through the canonical \texttt{METRICS\_INSERT} statement as a batched
write. Two policies keep the ingest path safe: a malformed payload returns a
distinct status from an unsupported content type, and a database write failure
returns success to the collector anyway, so the drop is counted and surfaced in a
stats endpoint while back-pressure stays away from the upstream. Only the
human-facing read API fails loudly, returning a 503 when the database is
unreachable, because an operator or downstream service needs to know the store is
down.

The \texttt{metrics} table is one of five base hypertables in the persistence
layer, each with a single authoritative writer; \cref{fig:ima:schema} lays them
out. It uses a wide-narrow shape, one row per reading carrying a
\texttt{metric\_name} and a \texttt{value}, so a new kind of metric needs no schema
change: the same write path that records \texttt{request\_latency\_ms} today records
any future signal tomorrow. On top of it sits \texttt{metrics\_1min}, a continuous
aggregate materialised every minute, which lets the decision engines and dashboards
ask minute-resolution questions without scanning the raw rows. The remaining tables,
\texttt{backend\_health}, \texttt{scaling\_events}, \texttt{forecasts}, and
\texttt{policy\_changes}, are each owned by the one service that produces them, the
same single-writer discipline applied to history rather than live telemetry.

\begin{figure}[tbp]
\centering
\includegraphics[width=\linewidth,height=0.88\textheight,keepaspectratio]{schema}
\caption[The persistence schema and its writers]{The five TimescaleDB hypertables plus the \texttt{metrics\_1min} continuous aggregate, each base table having exactly one authoritative writer.}
\label{fig:ima:schema}
\end{figure}

\cref{fig:ima:dataflow} traces the complete request-plus-telemetry path. The
request path (top) is synchronous and on the critical path; the telemetry path
(bottom) is asynchronous and lossy by design, so no part of the metric pipeline can
slow a client response.

\begin{figure}[H]
\centering
\begin{tikzpicture}[node distance=7mm and 14mm]
  \node[slnode] (client) {Client};
  \node[slnode,right=of client] (nginx) {NGINX\\load-balancer};
  \node[slnode,right=of nginx] (backend) {backend\\(served)};

  \node[slnode,below=10mm of nginx] (log) {access.log\\(JSON, shared vol.)};
  \node[slnode,below=8mm of log] (shipper) {lb-otel-shipper\\(tail $\to$ OTLP)};
  \node[slnode,right=of shipper] (otel) {otel-collector};
  \node[slnode,right=of otel] (tele) {telemetry\\(parse + insert)};
  \node[slnode,below=8mm of tele] (db) {TimescaleDB\\\texttt{metrics}};
  \node[slnode,left=of db] (dp) {decision plane\\(SQL reads)};

  % request path (synchronous)
  \draw[slflow] (client) -- node[font=\scriptsize,above]{req} (nginx);
  \draw[slflow] (nginx) -- node[font=\scriptsize,above]{proxy} (backend);
  \draw[slflow] (backend) to[out=200,in=-20] node[font=\scriptsize,below]{resp} (nginx);
  \draw[slflow] (nginx) to[out=160,in=20] (client);

  % telemetry path (asynchronous)
  \draw[slflow,dashed] (nginx) -- node[font=\scriptsize,right]{1 line/req} (log);
  \draw[slflow] (log) -- (shipper);
  \draw[slflow] (shipper) -- node[font=\scriptsize,above]{OTLP/HTTP} (otel);
  \draw[slflow] (otel) -- (tele);
  \draw[slflow] (tele) -- node[font=\scriptsize,right]{METRICS\_INSERT} (db);
  \draw[slflow] (db) -- (dp);
\end{tikzpicture}
\caption[The request and telemetry dataflow]{The synchronous request path and the asynchronous, lossy-by-design telemetry path, with the decision plane reading rows from TimescaleDB rather than live traffic.}
\label{fig:ima:dataflow}
\end{figure}

Read along the time axis, the same path falls into three regions. The data-plane
region is the only one a client waits on, and it ends not by calling the
telemetry pipeline but by appending a single line to the shared access log.
Everything downstream of that line, the tail, the OTLP hops, and the database
write, runs in its own time, while the control-plane engine cycle runs on a
separate clock again and reaches NGINX only through the sidecar reload. This
separation keeps the per-request overhead small and bounded regardless of how
slow the rest of the system happens to be.

One subtle point in the pipeline is backend identity. NGINX logs the upstream
\emph{address} (\texttt{\$upstream\_addr}), which resolves to an \texttt{ip:port}
rather than a container name; the shipper carries that address as the
per-data-point \texttt{instance} attribute, and the telemetry parser writes it into
\texttt{metrics.instance}. The decision plane keys on that field to decide
\emph{which} backend is misbehaving, so a later refinement added a reverse
resolution from backend IP to container name, matching the identity flowing through
the metrics to the one the sidecar and operator UI use. Without it, an anomaly
verdict and a routing weight could refer to the same backend under two different
names.

\subsection{Two observability paths}
\label{sec:ima:twopaths}

The pipeline above is only half of what the deployment observes. The design chapter
argued for keeping per-request telemetry and per-process health on separate,
vendor-neutral standards (\cref{dsa:sec:standards-telemetry}); the implementation
realises that split as two distinct paths that run side by side and never share a
transport, contrasted in \cref{tab:ima:observability}. The first is the high-volume
workload telemetry just described: one record per request, carried over OTLP into
TimescaleDB, where the decision engines read it. The second is the low-volume
service own-metrics path: each service exposes a Prometheus \texttt{/metrics}
endpoint that Prometheus scrapes directly, surfacing liveness, per-cycle timing, and
the decision counters that make the engines auditable, for instance
\texttt{anomaly\_detector\_isolate\_total} and \texttt{autoscaler\_scale\_total}.
Keeping the two apart means a flood of request telemetry cannot crowd out the signal
that tells an operator whether the engines themselves are alive.

\begin{table}[H]
\centering
\caption[The two observability paths side by side]{The two observability paths that run side by side: workload telemetry and service own-metrics.}
\label{tab:ima:observability}
\small
\begin{tabularx}{\linewidth}{@{}l Y Y@{}}
\toprule
 & \textbf{Workload telemetry} & \textbf{Service own-metrics} \\
\midrule
Question     & What the \emph{traffic} is doing & Whether the \emph{services} are healthy \\
Granularity  & One record per request           & One scrape per service per interval \\
Standard     & OpenTelemetry / OTLP \cite{opentelemetry} & Prometheus exposition \cite{prometheus} \\
Transport    & Log shipper, collector, telemetry & Direct \texttt{/metrics} scrape \\
Store        & TimescaleDB \texttt{metrics} hypertable & Prometheus time-series DB \\
Read by      & The decision engines (SQL) & Operators, via Grafana \cite{grafana} \\
Example signals & \texttt{request\_latency\_ms}, \texttt{error\_rate} & \texttt{<svc>\_up}, \texttt{<svc>\_cycle\_total}, \texttt{autoscaler\_scale\_total} \\
\bottomrule
\end{tabularx}
\end{table}

\section{Control Plane: Policy Management and Audit}
\label{sec:ima:policy}

The control plane begins at the policy manager, the single service that owns the
operating policy. It is the only writer of \texttt{config/policy.yaml}, the only
publisher on the \texttt{smartload.policy} channel, and the owner of the policy
audit trail. Every adaptive behaviour in the system, whether the autoscaler's
latitude, the latency multiplier that defines ``degraded'', or the confidence
threshold the routing engine must clear before it can act, is a field in this one
policy document, changeable only through the policy manager.

\subsection{Reading, validating, and writing policy}
\label{sec:ima:policyrw}

The policy manager exposes two endpoints: \texttt{GET /api/v1/policy} returns the
current policy, and \texttt{POST /api/v1/policy} commits a change. The POST is a
multi-step flow whose order is the heart of the service's correctness. Validation
comes first, written as pure functions with no I/O dependency so the full rule set
is unit-testable on its own. It runs in three passes: unknown keys in the request
body are rejected outright; each user-settable field is shape-checked on its own (an
integer field rejects a boolean, because in Python a boolean \emph{is} an integer
and would otherwise slip through); and the merged result is checked for cross-field
rules such as \texttt{min\_backends} $\leq$ \texttt{max\_backends}. Only after
validation succeeds and a non-empty diff is computed does the service write.

The write is atomic. The new policy is serialised to a temporary file in the same
directory and renamed over the canonical path, a rename atomic on both POSIX and
Windows. Docker single-file bind mounts add a subtlety: they mount the inode, not
the directory entry, and so reject the rename with \texttt{EBUSY}. The service
handles this with a documented fallback to a single buffered overwrite, bounding the
partial-read window to one operating-system write. The on-disk YAML is the source of
truth: the two steps that follow, the audit record and the bus publish, are both
best-effort. If both fail, the operator's commit has still landed, and subscribers
pick up the change on their next poll.

\subsection{Per-field audit and the policy envelope}
\label{sec:ima:audit}

Each committed change writes one row per changed field into the
\texttt{policy\_changes} hypertable: the field name, its old and new values
(JSON-encoded so the column holds any type), the policy version, the timestamp, and
the actor taken from an \texttt{X-Actor} header. One row per field, rather than one
per commit, lets the audit answer field-level questions directly: ``who last changed
\texttt{safe\_mode}, and to what?'' becomes a single filtered query rather than a
diff over snapshots.

After the write and audit, the policy manager publishes the full canonical policy
snapshot, every field rather than a delta, on the \texttt{smartload.policy}
Redis~\cite{carlson2013redis} channel, wrapped in the standard event envelope. The
snapshot carries a monotonically increasing \texttt{policy\_version} and a
\texttt{changed\_fields} list, so each subscriber can log what changed without
diffing snapshots. Sending the whole snapshot means a subscriber that missed an
earlier publish is still brought fully up to date by the next one; the version
number lets the autoscaler and sidecar reject a stale re-publish that would
otherwise rewind a live policy.

The audit log the operator sees is the union of two streams that deliberately stay
in separate tables owned by separate services: \texttt{policy\_changes} on the
policy manager and \texttt{scaling\_events} on the autoscaler. The owner of a write
is the owner of its read. They are unified only at the surface, where the SDK and
operator UI present one logical audit view, and never consolidated server-side,
keeping each service authoritative over its own history.

\section{The lb-sidecar: Applying Decisions to NGINX}
\label{sec:ima:sidecar}

Decisions do nothing until something applies them to NGINX. That something is the
\texttt{lb-sidecar}, a Python service running alongside the load balancer that
translates control-plane events into live NGINX configuration, bridging the
abstract intent of the decision plane and the concrete upstream block NGINX serves
from. The sidecar subscribes to four Redis channels and, on every message,
atomically rewrites an \texttt{upstream.conf} on a shared volume and sends NGINX a
reload signal (\texttt{nginx -s reload}). The reload swaps the upstream block live
without dropping in-flight connections.

\subsection{The four-channel dispatch loop}
\label{sec:ima:dispatch}

Each channel maps to one pure handler, and the dispatch loop is a single branch
on the channel name. \cref{tab:ima:channels} lists the four. Before every handler
runs, the sidecar re-queries the live container set from the Docker daemon rather
than trusting the count carried in the event. This makes the loop closed: every
rewrite reflects the containers that actually exist, not the publisher's intent,
so a backend whose provisioning failed at the Docker layer never appears in the
upstream block.

\begin{table}[H]
\centering
\caption[The four control channels the lb-sidecar subscribes to]{The four control channels the lb-sidecar subscribes to.}
\label{tab:ima:channels}
\small
\begin{tabularx}{\linewidth}{@{}l l Y@{}}
\toprule
\textbf{Channel} & \textbf{Publisher} & \textbf{Effect on NGINX} \\
\midrule
\texttt{smartload.routing} & rl-engine & If active and above the confidence threshold, convert backend rankings to integer upstream weights and rewrite \\
\texttt{smartload.anomaly} & anomaly-detector & Render excluded backends as \texttt{server $\cdots$ down;} so NGINX skips them without removing them from the pool \\
\texttt{smartload.policy}  & policy-manager & Hot-reload runtime knobs; \texttt{safe\_mode} resets all weights to equal while preserving exclusions \\
\texttt{smartload.scale}   & autoscaler & Re-query the live pool and regenerate an equal-weight upstream map after a scale action \\
\bottomrule
\end{tabularx}
\end{table}

A message moves through the loop in four steps: an envelope arrives on one of the
four channels; the sidecar re-queries the live backend set; the matching handler
computes a new weight map (or exclusion set); and the NGINX file adapter writes
\texttt{upstream.conf} and signals the reload. All four handlers converge on that
single adapter, which short-circuits when the computed map is identical to the one
already on disk, so no reload occurs. The \texttt{smartload.scale} channel makes
the closed scaling path explicit, turning an autoscaler action into a regenerated
equal-weight map and a live reload.

\subsection{Exclusion preservation and startup hydration}
\label{sec:ima:exclusion}

Two refinements give the sidecar its resilience. First, \emph{anomaly exclusions
survive partial routing updates}. A routing recommendation may rank only a subset of
backends; if the handler simply wrote those, a previously excluded backend could be
silently restored. Instead, each routing rewrite starts from a floor-weight baseline
over the full known pool and overlays the routing engine's values, so an excluded
backend always still appears in the file, rendered as \texttt{down;} by the adapter,
no matter which handler triggered the rewrite. Anomaly-excluded backends are running
containers that happen to be unhealthy, so they stay in the live pool and the
adapter's exclusion path renders them \texttt{down} whatever the cause of the
rewrite.

Second, \emph{exclusions are hydrated from the database on startup}. On boot, the
sidecar reads the latest health verdict per backend from the \texttt{backend\_health}
table and pre-populates its excluded set, so an exclusion persists across a Redis
disconnect or sidecar restart rather than silently lapsing. This step is non-fatal:
if the database is unreachable, the sidecar starts with an empty exclusion set and
rebuilds it from the next round of anomaly events.

The dispatch path also encodes two safety guards that bound the loop. If the
Docker query returns zero backends, whether from a daemon outage or a network
partition, the handler refuses to write an empty upstream block and leaves the
previous configuration intact, so NGINX keeps serving the last known good pool
rather than 502-ing every request. A quorum guard refuses an exclusion that would
empty the pool entirely: an \texttt{unhealthy} event that would exclude the last
active backend is turned into a no-op, because an empty upstream would 502 the
whole pool and feed back as still more exclusions.

\section{Operator Interface}
\label{sec:ima:ui}

The operator UI is the human face of the control plane. Following the design lock,
it is a \emph{transparency and override} layer, not an administration panel: it lets
an operator see what the decision plane is doing, propose a policy change with a diff
preview, browse the audit history, and trigger manual interventions, but it never
holds authority the API does not. It is built as a Flask backend-for-frontend (BFF)
serving a React front end compiled by Vite, and deploys as a single multi-stage
container in which the Node build stage produces static assets that the Python
runtime stage serves alongside its API.

The BFF plays three roles. It fans out \texttt{/health} across the whole stack in
parallel for the Flightdeck view; it proxies the policy, actions, and audit APIs
so the browser talks to a single origin; and it runs a Redis subscriber thread
that feeds a live event stream. The browser reaches only the BFF, which reaches
the rest of the stack on its behalf, so there is no direct browser-to-service
edge and authentication, CORS, and rate limiting are consolidated in one place.
The React app organises ten views into three task groups, surveyed in
\cref{sec:ima:pages}.

\subsection{The console views}
\label{sec:ima:pages}

Programmatic integrators drive SmartLoad through the SDK and webhooks, while the web
views give an operator a live, inspectable picture of the closed loop and a small
set of safe manual controls. The views are organised into three task groups, and one
principle governs them all: the data stays honest. Every indicator, delta,
sparkline, confidence band, and forecast is computed from real telemetry rather than
fabricated or frozen numbers. A \emph{Demonstration / Live} indicator shows which
mode is active and switches to live data the moment a backend is reachable, and
every panel carries explicit loading, empty, and error states so a missing source is
never disguised as a result.

\textbf{Overview.} \emph{Flightdeck} is the flagship view: a closed-loop overview
pulling throughput, tail latency, service-level attainment, and active incidents
into a single glance, alongside the current forecast and routing-and-scaling state.
\emph{System} draws the whole-system topology of the roughly a dozen services and
the data flow between them, including the observability pipeline, so the architecture
is legible from inside the console.

\textbf{Operate.} Five views expose the running planes. \emph{Pulse} shows
per-backend vitals, including the CPU and memory now collected per container.
\emph{Foresight} surfaces the demand forecaster and its scale-ahead signal.
\emph{Verdicts} is the anomaly-detector feed, carrying the evidence behind each
health verdict rather than a bare flag. \emph{Traffic} renders the load balancer
and sidecar view: request distribution, the current upstream weights, and any
exclusions. \emph{Capacity} is the autoscaler view: pool size against target, the
scale-ahead trajectory, and the scaling audit and heartbeat.

\textbf{Decide.} Three views hold the decision and control surfaces.
\emph{Helmsman} is the routing-engine view: the learned policy's shadow
evaluation, its proposed-versus-applied share, and its deploy-time promotion
readiness. Consistent with the safety model, routing is shadow by default and
promotion is a deploy-time decision, so the console deliberately offers no live
``promote'' control. \emph{Controls} is the operating surface: the policy editor
with a live field-level diff preview, the \texttt{safe\_mode} kill switch, and the
manual scale, isolate, and weight overrides, each written to the audit trail.
\emph{Ledger} is the immutable audit trail across policy, scaling, and isolation
events, with filters and CSV export for offline analysis.

The console ships a complete light and dark theme with a persisted toggle, a
responsive layout that collapses the sidebar into a mobile drawer, and accessibility
throughout: visible focus, focus traps on dialogs, reduced-motion support, and
contrast that meets the WCAG guidelines.

\subsection{The live stream and the status aggregator}
\label{sec:ima:sse}

The live engine views need sub-second updates as events arrive, too fine-grained for
polling. The BFF therefore runs a daemon Redis subscriber that consumes the four
event channels into a bounded per-channel ring buffer and exposes \texttt{GET
/api/ui/engines/stream} as a Server-Sent Events stream: new clients receive a replay
of recent events and then a live feed, with periodic heartbeat comments to keep the
connection open. The BFF runs threaded gunicorn workers, so a long-lived stream
occupies one thread rather than pinning a whole worker.

The Flightdeck view and any programmatic operator are served by a single aggregator,
\texttt{GET /api/v1/status}, that collapses the whole operational picture into one
document: every service's health, the active policy's headline fields, and the most
recent row from each audit stream. The fan-out is parallel with a two-second
per-service timeout, so one hung service can only delay the response by that bound
rather than serialising into an unbounded wait; an unreachable service shows up as
\texttt{down} in the per-service map, and an \texttt{overall} field rolls the whole
stack up into a single pill. The endpoint always returns 200, since callers route on
the \texttt{overall} field and per-service status rather than the HTTP code, so a
status read never becomes a point of failure.

\subsection{The Python SDK}
\label{sec:ima:sdk}

The same \texttt{/api/v1} surface the UI consumes is also a Python SDK, so an
integrator can drive SmartLoad programmatically. The SDK is organised as
sub-clients, one per concern. Five are fully implemented: \texttt{PolicyClient}
(read, set with diff, subscribe to changes), \texttt{StatusClient} (the consolidated
status read), \texttt{ActionsClient} (manual scale and isolate),
\texttt{AuditClient} (both history streams), and \texttt{EnginesClient} (engine
snapshots and a live subscription). Two further sub-clients, \texttt{MetricsClient}
and \texttt{WebhooksClient}, are deliberately stubbed: they raise
\texttt{NotImplementedError} and are planned for Phase~2, because raw metric access
and outbound webhook delivery build on the multi-tenancy and API-key work scheduled
for that phase. The deferral is a scope choice declared in the code rather than left
implicit.

\section{Iterative Development Process}
\label{sec:ima:iterative}

None of the modules above was built in one pass. We grew the control plane and
operator UI through a \emph{vertical slice} discipline: each feature was driven end
to end, through every product surface at once, before the next began. A slice was
done not when its service code compiled but when the same behaviour existed as an
HTTP endpoint in the canonical OpenAPI~\cite{openapi} specification, as a method on
the Python SDK, as a runnable scenario script, as an end-to-end test against the
live stack, and as a page or control in the operator UI. This kept the twelve
services coherent: a feature could not drift into one surface without the others,
because the slice was complete only when all surfaces agreed.

The slices landed in dependency order, which is itself the story of how the control
plane came together. \emph{Policy management} came first (\cref{sec:ima:policy}):
without a policy surface there was nothing for later features to read or change, so
the first slice delivered read, validate, diff, commit, publish, and audit across
HTTP, Redis, SDK, scenario, test, and UI, becoming the template every later slice
copied. \emph{Audit} followed, building on policy to give operators a unified
investigation view across both history streams without consolidating them
server-side. \emph{Manual actions} came next, reusing the audit pattern so that
operator-initiated scale and isolate land in the same history as automated
behaviour, set apart only by their \texttt{manual:} prefix. \emph{Live Engines}
added the real-time decision feed, requiring the BFF to grow a Redis subscriber and
an SSE transport. The \emph{consolidated status} slice then collapsed the seven-call
health-polling burst into one aggregated read. Finally, the \emph{lb-sidecar}
(\cref{sec:ima:sidecar}) closed the actuation loop, turning the decisions that all
the earlier surfaces could observe into live changes to the NGINX upstream.

We refined each slice repeatedly after its first landing. The sidecar alone gained
exclusion preservation, confidence gating, startup hydration, a scale channel, and a
quorum guard across several follow-on iterations, each a small, separately tested
increment rather than a rewrite. That cadence, shipping the thinnest end-to-end
slice and then hardening it in place under a running stack, is the development
process that produced the deployed system this chapter has described. The
decision-plane services and the models they run, which the slice discipline kept
observable from their first day through the live engine views, are the subject of
Part~B.
